% ================================================================
%  chapters/parts/ch10-b.tex
%  第 10 章  多相催化反应动力学（分片 b）
%  源：S2（下册）第 10 章，书内 p.769—792 / PDF p.258—281
%  覆盖：10.3 动力学模型的建立与检验；10.4 多相催化中的传递过程
%  说明：分片 a 写 10.1—10.2，分片 c 写 10.5—10.6。本分片不含 \chapterstart。
%        所有公式按「OCR 上下文 + 学科标准形式」重建，未对照原页面图像；
%        存疑处一律就地标 \tocheck{}。
% ================================================================

\subsection{动力学模型的建立与检验}
\label{sec:ch10-model}

前两节给出了速率方程的各种形式和判断速率控制步骤的方法。但从一组实验数据到一个可以
外推的速率方程，中间还有三步工作：定候选模型族、估参数、判定「这个模型能不能用」。
本节做这三步，并讲清一句常被含糊过去的话——\textbf{拟合优度高不等于模型正确}。

\begin{probout}
（1）面对同一组「速率—组成」数据，幂式、L-H 式、非均匀表面的 Temkin 式在数据窗口内
可能几乎重合，凭什么选一个？（2）参数估出来以后，怎么判断它\textbf{物理上合理}，而不
只是数学上的一组数字？（3）残差图该看什么——「残差小」与「残差随机」是两回事；
（4）什么条件下两个模型\textbf{原则上不可能}被这组数据区分（可辨识性），此时继续比较
$R^2$ 等于在比较噪声。本节还要交代一件更少被提及的事：\textbf{估计参数所用的数据窗口
本身决定了参数的含义}——同一个模型、同一套数据，只换窗口，表观级数就会变。
\end{probout}

\subsubsection{模型族：从幂式到非均匀表面}

催化剂动力学模型分三个层次，混用层次是很多错误的起点。\textbf{幂式}
$r=kp^n$ 是\textbf{纯经验的插值式}：$n$ 不是机理量，它只保证在实测窗口内可用；
\textbf{双曲线式（L-H 型）}由吸附平衡加一个速率控制步骤推出，参数有明确含义
（吸附平衡常数 $K$、表面速率常数 $k_S$），可外推到未测的组成；\textbf{非均匀表面模型}
再往前走一步，不假定所有活性位等价，而是给出一套吸附势分布。前两节的 L-H 式与
M-vK 式属于第二层，下面补上第三层。

\begin{principle}
\textbf{非均匀表面（Temkin）模型的基本假定}\src{S2 §10.3，书内 p.769 / PDF p.258}：
（1）活性位的标准吸附势 $A_0$ 有连续分布函数
$\dd S'=a\exp(-\gamma A_0/RT)\,\dd(A_0/RT)$，$\gamma$ 为分布特征参数，$a$ 为归一化因子，
$\dd S'$ 是标准吸附势落在 $A_0$ 与 $A_0+\dd A_0$ 之间的活性位数目；
（2）速率常数与平衡常数之间有 Br\"onsted 型关系 $k_i=cK_i^{\alpha}$，$\alpha$ 为
Temkin 系数，取值 $0$—$1$，通常为 $0.5$；
（3）化学上类似的两个步骤（吸附与脱附）的变换系数相同；
（4）覆盖度 $t$ 的上限是全部活性位被占，下限是全部活性位空着。
\end{principle}

由这四条假定，可对两步反应 $\mathrm{S_1+A_1\rightleftharpoons B_1+S_2}$、
$\mathrm{S_2+A_2\rightleftharpoons B_2+S_1}$ 导出非均匀表面上的速率式
\src{S2 §10.3，书内 p.769 / PDF p.258}：

\begin{equation}
v_t=\frac{v}{[L]}
=\tau\,
\frac{k_1^{\mathrm{o}}k_2^{\mathrm{o}}[A_1][A_2]-k_{-1}^{\mathrm{o}}k_{-2}^{\mathrm{o}}[B_1][B_2]}
{\bigl(k_1^{\mathrm{o}}[A_1]+k_{-2}^{\mathrm{o}}[B_2]\bigr)^{m}
 \bigl(k_{-1}^{\mathrm{o}}[B_1]+k_2^{\mathrm{o}}[A_2]\bigr)^{1-m}}
\label{eq:temkin}
\end{equation}

其中 $v/[L]$ 是单位活性位上的速率，$k_i^{\mathrm{o}}$ 为 $A_0=0$ 处的速率常数；
分布因子 $\tau$ 与指数 $m$ 由式 (10-99) 的后两行给出。
\tocheck{未对照原页面图像：式 (10-99) 中 $\tau$ 与 $m$ 的确切定义——OCR 只剩
$\pi/\sin(\pi m)$ 与 $\exp(\gamma f)-1$ 的残迹，$\exp(\gamma f)-1$ 在分子还是分母、
$m$ 与 $a$、$\gamma$、$f$ 的组合方式都无法从文字层确认；引用前请核对原页面。}

式~\eqref{eq:temkin} 真正值得看的是\textbf{结构}：分母是两个「吸附项之和」分别取 $m$ 与
$1-m$ 次幂。当某个物种的吸附项占优时，表观级数自动落在 $0$ 与 $1$ 之间连续变化——
这正是非均匀表面模型与 L-H 式在形式上的最大差别（L-H 式的幂次是整数，级数只能取
$1$、$0$、$-1$ 这些离散值）。原书指出，该式除由两步反应导出外，还得到铁基催化剂上
氨合成的总包动力学数据、氮气吸附等温线形式和氮吸附速率数据的支持。
\src{S2 §10.3，书内 p.769 / PDF p.258}

\begin{table}[htbp]
\centering
\caption{同一组数据上的四类候选速率方程：参数、可辨识条件与失效边界。
选取时应先看第 4 列——数据窗口是否覆盖了能把模型区分开的那段条件。
\src{据 S2 §10.2.2—§10.2.3（书内 p.753—756 / PDF p.242—245）、§10.2.4
（书内 p.757—759 / PDF p.246—248）、§10.3.5（书内 p.768—769 / PDF p.257—258）整理。}}
\label{tab:ch10-modelfam}
\begin{tabularx}{\textwidth}{@{}l l l X@{}}
\toprule
类型 & 速率方程（单组分，$p$ 为分压） & 参数 & 可辨识条件与失效边界 \\
\midrule
幂式 & $r=kp^{n}$ & $k,\ n$ & $n$ 是窗口的产物而非机理量（本节例：$n=0.50\to0.20$）；
只可在实测窗口内插值，不可外推到极值区 \\
L-H 单站 & $r=kKp/(1+Kp)^{2}$ & $k,\ K$ & 必须覆盖 $p\approx1/K$ 的极值区；$k$ 与 $K$
强相关时须由独立吸附数据定 $K$，否则两者只能定出乘积 \\
L-H 解离 & $r=k\sqrt{Kp}/(1+\sqrt{Kp})^{2}$ & $k,\ K$ & 窗口窄时与单站式\textbf{不可区分}
（$p\ll1/K$ 时两式都退化为一级） \\
非均匀表面 & 式~\eqref{eq:temkin} & $k_i^{\mathrm{o}},\ \gamma,\ m$ & 参数多，需要覆盖
吸附项主次互换的条件区；单一温度的数据无法定出分布参数 $\gamma$ \\
\bottomrule
\end{tabularx}
\end{table}

\subsubsection{参数估计：数据窗口决定参数的含义}

\textbf{先定反应器，再定转化率范围。}求取动力学参数的反应器按转化率分为两类
\src{S2 §10.5.1，书内 p.779—781 / PDF p.268—270}：\textbf{积分反应器}装填足量催化剂，
进出口组成差别显著，无法用某个平均值代表全床组成，得到的是\textbf{转化率对空时的积分
数据}，分析要求不高，但必须对模型作积分后才能与数据比较；\textbf{微分反应器}的转化率
很低，进出口浓度差小到可用平均值代表全床组成，于是 $\Delta c/\Delta t$ 可近似为
$\dd c/\dd t$ 并等于反应速率 $r$，直接得到「$r$ 对平均组成」的\textbf{微分数据}，
代价是分析精度要求很高，而且数据往往是初速，难以配出复杂反应在高转化下的组成
（对策是在微分反应器前串接一个积分转化器提供进料）。
单程流通的管式微分反应器中，\textbf{转化率应在 $5\%$ 以下（个别允许到 $10\%$）}。

\begin{measure}
\textbf{估参数之前要固定的四件事。}（1）\textbf{恒温}：温度对速率是指数的，任何温度
梯度都直接变成速率误差。（2）\textbf{接触时间确切}：接触时间不确定，模型就无法与数据
对齐。（3）\textbf{转化率范围}：微分数据要求 $X<5\%$（个别 $\le10\%$）；
\textbf{转化率超过约 $15\%$ 的数据不能用来求反应级数}——此时沿床层的组成变化太大，
「平均组成」不再代表反应发生时的组成，级数被系统性地推离真值（见下方误判块）。
（4）\textbf{独立条件数}：独立自变量（组成、温度）的个数必须多于待估参数个数；
温度点少于 $3$ 个时 Arrhenius 参数与指前因子无法分开。
\src{S2 §10.5.1，书内 p.779—781 / PDF p.268—270}
\end{measure}

\textbf{拟合的权重问题。}把 $r=kp^n$ 取对数化成直线再作最小二乘，是最常见的做法，但它
\textbf{把「对 $r$ 的等权拟合」偷换成了「对 $\ln r$ 的等权拟合」}：低速率点的相对误差被
放大，高速率点被压缩。差别在窗口宽时可达可观量级。稳妥做法是：用对数直线只求\textbf{初值}，
再用非线性最小二乘直接拟合 $r$（或拟合相对残差 $\delta_i=(r_i^{\mathrm{cal}}-r_i)/r_i$，
使各点权重与其精度相称）。

\textbf{参数相关性是估计阶段的主要风险。}在 $r=kKp/(1+Kp)^2$ 里，只要 $Kp\ll1$，
速率就退化为 $r\approx kKp$：$k$ 与 $K$ 只能定出乘积 $kK$，各自的值由数据的噪声决定。
表现是：$k$ 与 $K$ 的相关系数接近 $\pm1$、置信区间极宽、换一批数据两者一起大幅漂移。
诊断与处置：算参数的相关矩阵，$|\mathrm{corr}|>0.95$ 时不要把 $k$ 与 $K$ 分开解释；
用独立信息（吸附量、程序升温脱附得到的覆盖度、瞬态实验）钉住其中一个参数。

\begin{dataflow}
\textbf{从原始数据到速率方程（本例为微分数据）。}
（1）核对物料衡算：进出口流量、组成、催化剂质量与接触时间的闭合误差；
（2）换算基准：把分析得到的摩尔数变化换成\textbf{单位质量}速率 $r_m$，需要与判据比较时
再乘颗粒密度换成单位体积速率 $r_v=\rho_pr_m$（详见下一节）；
（3）作图与选窗口：先画 $r$ 对 $p$ 的全窗口图，确认有没有极值、有没有产物阻抑或传质
控制的迹象，再决定用哪一段窗口估参数；
（4）求初值：对数直线或文献值；
（5）非线性最小二乘，输出参数、参数相关系数矩阵、残差序列；
（6）\textbf{判据检查}：残差是否随机、参数是否落在物理允许范围、窗口外能否预测
（见下小节）。
\end{dataflow}

换窗口会怎样，可以用本节的教学设定数据直接量出来：数据按 L-H 模型
（$k=2.00$、$K=0.50$\,\unit{kPa^{-1}}）生成并叠加 $\pm1.5\%$ 的相对噪声，共 $9$ 个点，
$p=0.5$—$3.0$\,\unit{kPa}。用同一个幂式分别拟合三个窗口，结果见表~\ref{tab:ch10-window}。

\begin{table}[htbp]
\centering
\caption{同一套教学设定数据、同一个幂式模型，只换拟合窗口：幂式指数 $n$ 从 $0.50$ 漂到
$0.20$（变 $2.5$ 倍），而 L-H 的吸附平衡常数 $K$ 在三个窗口下都停在 $0.51$—$0.52$
（生成值 $0.50$）。\textbf{参数的窗口稳定性比 $R^2$ 更能判断模型是否可信。}
\src{表中数值为本改编本按 L-H 模型生成的教学设定数据，非源书实验值；拟合窗口与参数
可辨识性的讨论据 S2 §10.3.2—§10.3.3，书内 p.760—765 / PDF p.249—254 整理。}}
\label{tab:ch10-window}
\begin{tabularx}{\textwidth}{@{}l c c c c@{}}
\toprule
拟合窗口 $p$/\unit{kPa} & 点数 & 幂式指数 $n$ & 幂式 $R^{2}$ & L-H 的 $K$/\unit{kPa^{-1}} \\
\midrule
$0.5$—$0.9$ & 3 & $0.496$ & $0.995$ & $0.510$ \\
$0.5$—$1.5$ & 5 & $0.371$ & $0.963$ & $0.520$ \\
$0.5$—$3.0$ & 9 & $0.201$ & $0.723$ & $0.510$ \\
\bottomrule
\end{tabularx}
\end{table}

窄窗口内幂式的 $R^2$ 高达 $0.995$，残差只有 $\pm1\%$：\textbf{只看 $R^2$，两个模型在这
三个点上无法区分}。但同一模型换到宽窗口，指数就从 $0.50$ 掉到 $0.20$；而 L-H 的 $K$
不随窗口移动。\textbf{一个参数如果随拟合窗口漂移，它就不是机理量。}
\pedset{表~\ref{tab:ch10-window} 与本节全部残差、$R^2$ 数值均为按 L-H 模型生成的
教学设定数据，非原书实验值。}

\subsubsection{模型判别与残差分析}

残差不是「拟合剩下的垃圾」，它是模型与数据之间唯一的结构性信息。判别的顺序是：
\textbf{先看残差的结构，再看参数的物理合理性，最后才看拟合优度}。本节用同一套数据的
两个拟合来演示。

\begin{figure}[htbp]
\centering
\begin{tikzpicture}
\begin{axis}[
  width=\onecolwidth,
  height=5.4cm,
  xlabel={分压 $p$ / \unit{kPa}},
  ylabel={相对残差 / \%},
  xmin=0.4, xmax=3.2, ymin=-15, ymax=15,
  xtick={0.5,1.0,1.5,2.0,2.5,3.0},
  ytick={-15,-10,-5,0,5,10,15},
  legend style={at={(0.5,1.02)}, anchor=south, legend columns=3, draw=none, fill=none},
]
\addplot[mainline, mark=*, mark size=1.5pt] coordinates {
  (0.50,12.11) (0.70,-0.12) (0.90,-5.57) (1.20,-6.37) (1.50,-7.70)
  (1.90,-4.64) (2.30,-0.79) (2.70,5.54) (3.00,9.60)};
\addlegendentry{幂式拟合}
\addplot[thinline, dashed, mark=square*, mark size=1.5pt] coordinates {
  (0.50,-0.03) (0.70,-0.43) (0.90,-0.54) (1.20,1.54) (1.50,-0.30)
  (1.90,-0.12) (2.30,-0.70) (2.70,0.29) (3.00,-0.04)};
\addlegendentry{L-H 拟合}
\addplot[accentsoft, dotted, thick] coordinates {(0.4,0)(3.2,0)};
\addlegendentry{零残差}
\end{axis}
\end{tikzpicture}
\caption{同一组数据、两种模型的相对残差随分压的分布。读图看两件事：幂式的残差
\textbf{符号成块}（$+,-,-,-,-,-,-,+,+$，连成一条弓形），L-H 的残差\textbf{符号交替}且
落在 $\pm1.6\%$ 内。条件与局限：数据为按 L-H 模型生成并加 $\pm1.5\%$ 相对噪声的教学
设定数据（$9$ 点，$p=0.5$—$3.0$\,\unit{kPa}），模型计算结果，不代表任何具体催化剂。
\modelmark}
\label{fig:ch10-resid}
\end{figure}

\begin{obs}
读图~\ref{fig:ch10-resid}：幂式拟合的残差在 $p=0.5$\,\unit{kPa} 处为 $+12.1\%$，随后
连续 $7$ 个点为负（$-0.1\%$、$-5.6\%$、$-6.4\%$、$-7.7\%$、$-4.6\%$、$-0.8\%$ 等），
在 $p>2.3$\,\unit{kPa} 又转为正（$+5.5\%$、$+9.6\%$）；符号序列为
$+,-,-,-,-,-,-,+,+$。L-H 拟合的 $9$ 个残差全部落在 $\pm1.6\%$ 以内，符号交替
（$-,+,-,-,-,-,+,-$），看不出结构。两个模型在同一组数据上的 $R^2$ 分别为
$0.723$ 与 $0.997$。
\end{obs}

\begin{infer}
（1）连续 $7$ 个同号残差不可能是噪声造成的：噪声在各点之间独立，不会成块同号。
\textbf{成块的残差是函数形式缺失的特征}——幂式在整段窗口里只有一个恒定斜率，而真实
速率曲线的斜率在变，误差于是按「两端正、中间负」的方式积累。（2）因此判别模型不能
用「残差小」或 $R^2$，要用\textbf{残差是否随机、是否与自变量或未进入模型的变量相关}。
（3）本例中 $R^2$ 恰好把错误模型暴露了（$0.723$ 对 $0.997$），但这是数据窗口的运气：
把窗口缩到表~\ref{tab:ch10-window} 的第一行（三个点），同一个幂式的 $R^2$ 就是
$0.995$、残差 $\pm1\%$——此时 $R^2$ 已经失去判别力。（4）残差随机只说明模型与这组
数据\textbf{相容}，不说明机理正确；机理结论必须另找独立证据（见本小节末的互补与回查）。
\end{infer}

\begin{table}[htbp]
\centering
\caption{残差图的形态与它提示的病因。这张表是判别阶段的检查清单：先按形态定位病因，
再决定是改模型、改权重还是去排除传递过程。
\src{本表为本改编本归纳的残差诊断清单，非源书表格；模型判别的准则据
S2 §10.3.4，书内 p.766—767 / PDF p.255—256。}}
\label{tab:ch10-resid-diag}
\begin{tabularx}{\textwidth}{@{}l X X@{}}
\toprule
残差形态 & 提示 & 处置 \\
\midrule
符号成块的连续弯曲 & 函数形式缺失（模型阶次或机理形式不对） & 换机理型模型；
扩宽条件窗口后重看残差 \\
残差与某个\textbf{未进入模型}的自变量相关（粒径、线速、空速） &
传递过程混入表观动力学 & 先按下一节的判据排除外扩散与内扩散，再重新拟合 \\
残差绝对值随速率增大 & 误差模型不对（等权拟合给了高值区过大权重） &
改用相对残差或加权最小二乘 \\
少数点残差远大于其余 & 离群点或分析误差 & 回查原始记录；不要为了 $R^2$ 删点 \\
残差随机但参数置信区间极宽、相关矩阵近奇异 & 数据信息量不足或参数强相关 &
补独立条件（更宽的分压、更多温度点、独立吸附量数据） \\
残差随机、参数稳定、窗口外可预测 & 与数据相容 & 仍不等于机理正确，需独立实验回查 \\
\bottomrule
\end{tabularx}
\end{table}

\textbf{「拟合优度高」为什么不能当判据。}三条理由，任一条都足以否掉用 $R^2$ 排模型优劣：
（1）$R^2$ 对\textbf{系统性偏差不敏感}——只要数据窗口窄、噪声大，错误模型也能拿到
$0.995$（表~\ref{tab:ch10-window} 第一行）；（2）$R^2$ 从不检查参数是否\textbf{物理
合理}，而 $k<0$、$K<0$、活化能为负的拟合在数值上完全可能；（3）增加参数\textbf{总能}
提高 $R^2$，所以参数个数不同的模型之间比 $R^2$ 是\textbf{过拟合的邀请函}——判据应是
「自由度 $n-p$ 下的残差是否与噪声量级一致」，而不是 $R^2$ 的大小。

\textbf{模型能否区分，取决于数据窗口。}把两个模型各自外推到数据窗口之外，差别才显现
（图~\ref{fig:ch10-extrap}）：幂式在 $p=8$\,\unit{kPa} 给出 $0.635$，比窗口内的最大值还
高；L-H 给出 $0.319$，并已越过 $p\approx2$\,\unit{kPa} 的极大值开始下降。实验数据在
$p>2$\,\unit{kPa} 确实下降，所以\textbf{幂式的函数形式在物理上已被否掉}——尽管它在
$0.5$—$0.9$\,\unit{kPa} 内 $R^2=0.995$。

\begin{figure}[htbp]
\centering
\begin{tikzpicture}
\begin{axis}[
  width=\onecolwidth,
  height=5.6cm,
  xlabel={分压 $p$ / \unit{kPa}},
  ylabel={相对速率（任意单位）},
  xmin=0.3, xmax=8.0, ymin=0, ymax=0.72,
  legend style={at={(0.98,0.03)}, anchor=south east, draw=none, fill=none},
]
\addplot[only marks, mark=*, mark size=1.6pt, forget plot] coordinates {
  (0.50,0.3241) (0.70,0.3893) (0.90,0.4331) (1.20,0.4629) (1.50,0.4911)
  (1.90,0.4985) (2.30,0.4980) (2.70,0.4835) (3.00,0.4756)};
\addplot[mainline, domain=0.3:8.0, samples=80]{0.4178*x^0.2014};
\addlegendentry{幂式外推}
\addplot[thinline, dashed, domain=0.3:8.0, samples=120]{1.99*(0.514*x)/((1+0.514*x)^2)};
\addlegendentry{L-H 外推}
\addplot[accentsoft, dotted, thick, forget plot] coordinates {(3.0,0)(3.0,0.72)};
\end{axis}
\end{tikzpicture}
\caption{把两种拟合外推到数据窗口之外。读图看 $p>3$\,\unit{kPa}（竖虚线右侧）两线如何
分道：幂式继续单调上升，L-H 越过极大值后下降。条件与局限：圆点为教学设定数据
（窗口 $p\le3$\,\unit{kPa}），两条曲线为模型计算；纵轴为任意单位，只比较形状。
\modelmark}
\label{fig:ch10-extrap}
\end{figure}

\begin{pitfall}
\textbf{本小节要防的四个误判。}（1）\textbf{用 $R^2$ 判定模型优劣}——理由见上；
排模型优劣要用残差结构、参数稳定性与窗口外预测。（2）\textbf{用转化率过高的数据求反应
级数}——单程转化率超过约 $15\%$ 后，沿床层组成变化过大，「平均组成」不再代表反应发生
时的组成，级数被推离真值；这是积分反应器数据被误当微分数据使用的典型后果。
（3）\textbf{把强相关参数当成两个独立结论}——$k$ 与 $K$ 的相关性接近 $\pm1$ 时，报出
「活化能」和「吸附热」两个数其实是同一份信息的两种说法。（4）\textbf{过拟合}——
用 $5$ 个参数拟合 $6$ 个数据点，残差必然很小，\textbf{自由度 $n-p$ 必须先于 $R^2$
报告出来}。
\end{pitfall}

\begin{crosscheck}
\textbf{参数的物理合理性检查（不依赖拟合优度）。}（1）速率常数与吸附平衡常数恒正；
（2）吸附放热，由 van't Hoff 关系 $\dd\ln K/\dd(1/T)=-\Delta H_{\mathrm{ads}}/R>0$，
故 $K$ 必须随温度升高而\textbf{减小}——若拟合出的 $K$ 随温度升高而增大，说明这个
「$K$」在吸收别的效应（例如把传质或失活算进去了）；（3）表观活化能应落在合理区间
（约 $20$—$300$\,\unit{kJ.mol^{-1}}），且与反应类型相称；强内扩散控制下它会被压到
真值的一半（见下一节）；（4）反应级数应在 $-3$ 至 $+3$ 之间，且正负号与抑制/促进关系
一致。若这四条与文献或独立实验冲突，模型即使 $R^2=0.999$ 也不能用。另一条回查路径是
\textbf{换一种实验方式复测}：用瞬态/同位素方法直接测吸附态的量与寿命，检验由稳态数据
反推的 $K$ 是否成立（方法见本章后半部分）。
\end{crosscheck}

\subsection{传递过程的影响与判据}
\label{sec:ch10-transport}

上一节假定测到的是表面化学反应的速率。这个假定很少自动成立：多相催化反应总是把物质
传递、热量传递与表面反应纠缠在一起，得到的「速率」里混着设备参数。本节的任务是把这
三股过程拆开、给出各自的参数与判据，并把「已排除扩散影响」这句话变成可执行的实验。

\begin{probout}
（1）传递过程究竟改变了速率方程的什么——是速率的大小，还是活化能与级数？（2）外扩散
与内扩散各自用什么参数描述，这些参数怎么测、关联式在什么范围内有效？（3）三个判据
（Damk\"ohler 数、Weisz--Prater 数、Mears 判据）分别检验哪一个环节，它们的\textbf{前提}
是什么、什么情况下会失效？（4）最要紧的一件事：三个判据里的速率必须是\textbf{同一种
基准}——催化研究习惯用单位质量速率，而判据是按单位颗粒体积推出来的，差一个颗粒密度；
本节对每一式做 SI 量纲展开，并给出同一条判据因基准不同把临界粒径从 $3.65$\,\unit{mm}
误判成 $0.105$\,\unit{mm}（相差 $35$ 倍）的实例。
\end{probout}

\subsubsection{七个步骤、耦合的后果与速率基准}

多相催化反应过程包含 $7$ 个步骤\src{S2 §10.4，书内 p.769—770 / PDF p.258—259}：
反应组分向催化剂外表面的传递（外扩散）$\to$ 由外表面向内表面的传递（内扩散）$\to$
在活性中心吸附 $\to$ 表面反应 $\to$ 产物脱附 $\to$ 产物由内表面向外表面传递（内扩散）
$\to$ 产物由外表面向流体体相传递（外扩散）。其中\textbf{外扩散与表面反应、内扩散是
连串的，而内扩散与表面反应是平行的}——这一结构差别决定了后面判据的顺序：先排外扩散，
再排内扩散。

耦合的后果有三条，都是可观测的\src{S2 §10.4，书内 p.770 / PDF p.259}：
（1）若忽视浓度梯度、把体相浓度当作表面浓度，求得的动力学方程\textbf{活化能偏低、
反应级数向 $1$ 靠近}；（2）恒温下由于传质阻力的存在，\textbf{测得的总包速率低于化学
反应的本征速率}；（3）对有显著焓变化的反应还叠加温度梯度，\textbf{动力学方程的失真
更严重}——相间与粒内的温差都会改变速率常数，而且是指数性的。因此正确的解析思路是把
全程分解为表面反应、相间传质、粒内传质和传热四个过程，分别求参数再纳入反应器模型；
从催化剂角度看，这样分解还能指出薄弱环节究竟是内扩散、导热性，还是本身的活性。
\src{S2 §10.4，书内 p.770 / PDF p.259}

\begin{principle}
\textbf{速率基准的声明（全节通用，也是三个判据的共同前提）。}动力学数据常按
\textbf{单位质量}催化剂报出，而传递过程的物料衡算是按\textbf{单位颗粒体积}写的，
两者只差一个颗粒密度（含内孔的颗粒密度 $\rho_p$，不是骨架密度）：
\begin{equation}
r_v=\rho_p\,r_m,\qquad
[r_m]=\unit{mol.kg^{-1}.s^{-1}},\quad
[r_v]=\unit{mol.m^{-3}.s^{-1}},\quad
[\rho_p]=\unit{kg.m^{-3}}
\label{eq:ch10-basis}
\end{equation}
\textbf{量纲检查}：$\unit{kg.m^{-3}}\times\unit{mol.kg^{-1}.s^{-1}}
=\unit{mol.m^{-3}.s^{-1}}$，闭合。本节的每一式都标明所用基准；凡是与判据比较的速率，
一律先换成 $r_v$。换成 $r_v$ 之后\textbf{不要再乘一次 $\rho_p$}——这一点在判据那一小节
会用一个图直接量出后果。
\end{principle}

\subsubsection{外扩散与相间传热}

流体体相向固体外表面的传质速率与推动力 $(c_B-c_S)$ 成正比，与扩散阻力成反比；用传质
系数 $k_G$ 表示阻力的倒数，则 $R=k_G(c_B-c_S)$。传质系数与流动状态、颗粒几何和体系
物性有关，准确的理论预测需要解传递微分方程，实用上仍靠无因次准数与经验关联式。
\src{S2 §10.4.1，书内 p.770—771 / PDF p.259—260}

\begin{principle}
\textbf{相间传质与传热的关联式}\src{S2 §10.4.1，书内 p.770—772 / PDF p.259—261}。
传质因子与 Sherwood 数：
\begin{equation}
j_D=\frac{\rho k_G}{G}\,Sc^{2/3}=\frac{Sh}{Re\,Sc^{1/3}},\qquad
Sc=\frac{\mu}{\rho D},\quad Sh=\frac{k_Gd_p}{D},\quad Re=\frac{Gd_p}{\mu}
\label{eq:ch10-jd}
\end{equation}
式中 $\rho$ 为流体密度，$G$ 为质量流速，$d_p$ 为颗粒直径，$D$ 为分子扩散系数。固定床在
$Re=0.8$—$2130$、$Sc=0.6$—$1300$ 范围内（气体、液体都适用）：
\begin{equation}
j_D=\frac{0.725}{Re^{0.41}-0.15}
\label{eq:ch10-jd-fix}
\end{equation}
外扩散成为控速步骤时，扩散组分自流相输向固体外表面的速率即该组分的转化速率，由物料
衡算可得转化率 $X$ 与床层长度的关系：
\begin{equation}
\ln\frac{1}{1-X}=\frac{0.725\,a\,z}{(Re^{0.41}-0.15)\,Sc^{2/3}}
\label{eq:ch10-extconv}
\end{equation}
式中 $a$ 为单位体积床层颗粒的表面积，$z$ 为床层高度。$ad_p$ 乘积不随颗粒直径变化，
在孔隙率 $\varepsilon\approx0.4$ 的典型催化床层中 $ad_p\approx3.6$（球形颗粒
$a=6(1-\varepsilon)/d_p$）。
对流传热的对应关系为
\begin{equation}
Q=R(-\Delta H_r)=\alpha_s a\,(T_S-T_B),\qquad
j_H=\frac{\alpha_s}{Gc_p}\,Pr^{2/3}=\frac{Nu}{Re\,Pr^{1/3}},\qquad
Nu=\frac{\alpha_s d_p}{\lambda_f}
\label{eq:ch10-jh}
\end{equation}
固定床在同一条 $Re$、$Sc$ 范围内 $j_H=1.10/(Re^{0.41}-0.15)$，两式相除即
$j_D/j_H=0.725/1.10=0.66$。\textbf{这里要留意}：经典的 Chilton--Colburn 类比给出
$j_D=j_H$（比值为 $1$），因此 $0.66$ 不是普适常数，而是\textbf{固定床在
$Re=0.8$—$2130$ 内两个经验关联式之商}；用其他床层结构或超出该 $Re$ 范围时比值要重算。
代入热量衡算得相间温差
\begin{equation}
T_S-T_B=\frac{j_D}{j_H}\Bigl(\frac{Pr}{Sc}\Bigr)^{2/3}
\frac{-\Delta H_r}{\rho c_p}\,(c_B-c_S)
\;\approx\;0.66\,\lambda f
\label{eq:ch10-dT}
\end{equation}
其中对大多数气体 $Pr/Sc\approx1$，$\lambda=(-\Delta H_r)c_B/(\rho c_p)$ 为放热反应的
绝热温升，$f=(c_B-c_S)/c_B$ 为流-固体系的浓度梯度。
\end{principle}

\textbf{两个可以当场算的检查。}（1）把 $ad_p=3.6$、$Sc\approx1$ 代入
式~\eqref{eq:ch10-extconv}：$Re=100$ 时 $Re^{0.41}=6.61$，$j_D=0.112$，
$X=90\%$ 要求 $z/d_p=5.7$——\textbf{大约 $6$ 粒催化剂长的床层就够了}。反过来说，如果
一个反应需要 $50$—$100$ 层催化剂才能达到 $90\%$ 转化率，就可以认为它\textbf{不是外扩散
控速}，气流中的分压基本等于颗粒外表面的分压。这把「有没有外扩散限制」变成了一次
可算的比较。\src{S2 §10.4.1，书内 p.771 / PDF p.260}
（2）式~\eqref{eq:ch10-dT} 的数量级含义：$T_S-T_B$ 与反应热成正比，与浓度梯度成正比。
\textbf{即使相间传质限制不大（$f$ 小），只要 $|\Delta H_r|$ 足够大，相间温差仍可把速率
推得远高于无传热影响时的值}——因为速率对温度是指数的。这是「传质已排除」不能推出
「传热已排除」的原因。\src{S2 §10.4.1，书内 p.772 / PDF p.261}

\tocheck{未对照原页面图像：式 (10-100) 之后对 $c_B$ 与 $c_S$ 的文字说明在 OCR 中与
全书其余各处\textbf{相反}（OCR 作「$c_B$ 为颗粒外表面浓度、$c_S$ 为体相浓度」，而
式 (10-130) 明写「$c_B$ 为组分 $\mathrm{B_1}$ 在流体相的浓度」，式 (10-107) 的
$T_S$、$T_B$ 也是「表面/体相」）。本节按 $B$=体相、$S$=外表面书写，引用原式前请核对
页面。}

\begin{pitfall}
\textbf{外扩散环节的边界。}（1）式~\eqref{eq:ch10-jd-fix} 与
式~\eqref{eq:ch10-jh} 只在 $Re=0.8$—$2130$、$Sc=0.6$—$1300$ 内经验证，
\textbf{不能外推}到低线速或高黏体系（实验室微型反应器常在 $Re<1$）。（2）
$ad_p\approx3.6$ 只对孔隙率约 $0.4$ 的球形颗粒床成立；细粉、规整填料或大孔载体要重算。
（3）\textbf{把「传质已排除」当成「无温度梯度」}——强放热反应的热点可以出现在传质限制
很小的时候（见式~\eqref{eq:ch10-dT} 后的讨论）。（4）用改变流量去检验外扩散时，
必须\textbf{同步改变催化剂用量以保持接触时间不变}；只变流量而不同步变催化剂量，
测到的是接触时间的影响，不是线速的影响。
\end{pitfall}

\subsubsection{粒内传质：Thiele 模数与有效因子}

反应物从外表面进入颗粒内孔，按分子性质与孔结构分四种机制
\src{S2 §10.4.2，书内 p.772—773 / PDF p.261—262}：孔径远大于分子平均自由程时是
\textbf{分子（正常）扩散}（常压下平均自由程约 $10^{-7}$\,\unit{m} 量级，与压力成正比）；
孔径远小于平均自由程时，分子与孔壁碰撞的机会远高于分子间碰撞，是 \textbf{Knudsen 扩散}；
孔径接近分子尺寸时\textbf{构型扩散}（扩散系数与被吸附分子构型强相关，可差几个数量级，
顺-2-丁烯与反-2-丁烯在 A 型分子筛中的扩散系数相差 $200$ 倍以上）；吸附分子在表面
移动产生的输运是\textbf{表面扩散}。实际颗粒的孔道复杂，故引入\textbf{有效扩散系数}。

\begin{table}[htbp]
\centering
\caption{四种孔内扩散机制：系数、依赖关系与适用孔径。第 3 列是判据用得上的一条——
压力依赖不同，会推出不同的\textbf{表观反应级数}（见本节后文）。
\src{据 S2 §10.4.2，书内 p.772—775 / PDF p.261—264 归纳。}}
\label{tab:ch10-diffmech}
\begin{tabularx}{\textwidth}{@{}l X l X@{}}
\toprule
机制 & 有效系数 & 温度、压力依赖 & 适用条件 \\
\midrule
分子（正常）扩散 & $D_{ej}=\varepsilon_p D/\Gamma$ & $D\propto T^{3/2}/p$，
与孔径无关 & 孔径 $\gg$ 分子平均自由程（常压下约 $>0.1$\,\textmu m） \\
Knudsen 扩散 & $D_K=97\,r_p\sqrt{T/M}$，$D_{Ke}=\varepsilon_pD_K/\Gamma$ &
$\propto r_p\sqrt{T}$，与压力无关 & 孔径 $\ll$ 分子平均自由程 \\
过渡区 & $1/D_e=1/D_{Ke}+1/D_{ej}$ & 由两者中较小者控制 &
两种机制贡献相当；$D_e$ 与气体组成有关 \\
表面扩散 & $D_s=D_0\exp(-E_s/RT)$ & 活化能约为物理吸附活化能的一半 &
吸附层内输运，常温下 $D_s$ 约 $10^{-7}$—$10^{-5}$\,\unit{m.s^{-1}} \\
构型扩散 & 与分子构型强相关 & 随孔口尺寸与分子形状变化 &
孔径接近分子尺寸（分子筛）；数值可差几个数量级 \\
\bottomrule
\end{tabularx}
\end{table}

\tocheck{未对照原页面图像：表~\ref{tab:ch10-diffmech} 中 $D_s$ 的数值区间原书作
「$10^{7}$—$10^{5}$ m/s」，上标符号在 OCR 中丢失；按物理量级应为 $10^{-7}$—
$10^{-5}$，且单位是否应为 \unit{m^2.s^{-1}} 无法从文字层确认。}

\textbf{简化假定与浓度分布。}为定量解析，先考虑单一反应物并假定
\src{S2 §10.4.2，书内 p.773 / PDF p.262}：①恒温；②扩散遵从 Fick 法则，可用一个恒定的
总包有效扩散系数表征；③反应不可逆、无分子数变化、可用简单速率方程描述；④系统处于
稳态（扩散最深部浓度梯度为零）。颗粒取两种几何：半径 $R$ 的圆球，厚度 $L$ 的平板
（侧面封闭，仅正面接触反应物）。外表面浓度 $c_S$，一边以 $D_e$ 向内扩散、一边以
$k_SS_vc^m$ 转化，则离球心 $r$ 处、离板底 $l$ 处的浓度为
\src{S2 §10.4.2，书内 p.773—774 / PDF p.262—263}：

\begin{equation}
\frac{c}{c_S}=\frac{\sinh(\varphi_S r/R)}{(r/R)\sinh\varphi_S}\ \text{（球形）},\qquad
\frac{c}{c_S}=\frac{\cosh(\varphi_L l/L)}{\cosh\varphi_L}\ \text{（平板）}
\label{eq:ch10-profile}
\end{equation}

\begin{equation}
\varphi_S=R\sqrt{\frac{S_vk_S}{D_e}},\qquad
\varphi_L=L\sqrt{\frac{S_vk_S}{D_e}}
\label{eq:ch10-thiele}
\end{equation}

式中 $S_v$ 为单位颗粒体积的内表面积，$k_S$ 为单位内表面积的速率常数，
$\varphi_S$、$\varphi_L$ 为无因次的 \textbf{Thiele 模数}。
产物与反应物的浓度在颗粒内部都低于外表面，\textbf{内部的实际反应速率也低于外表面
条件下的速率}。稳态下总包速率与扩散速率相等，对圆球
\src{S2 §10.4.2，书内 p.774 / PDF p.263}：

\begin{equation}
r=4\pi R\,\varphi_SD_ec_S\Bigl(\frac{1}{\tanh\varphi_S}-\frac{1}{\varphi_S}\Bigr),
\qquad
r_0=\frac{4}{3}\pi R^{3}S_vk_Sc_S
\label{eq:ch10-r-r0}
\end{equation}

式~\eqref{eq:ch10-r-r0} 中 $r$ 为有内扩散时的颗粒速率，$r_0$ 为无扩散限制（内部浓度
处处等于 $c_S$）时的速率。两者之比就是\textbf{有效因子} $\eta$（催化剂体积利用率）：

\begin{equation}
\eta=\frac{r}{r_0}=\frac{3}{\varphi_S}
\Bigl(\frac{1}{\tanh\varphi_S}-\frac{1}{\varphi_S}\Bigr)\ \text{（球形）},\qquad
\eta=\frac{\tanh\varphi_L}{\varphi_L}\ \text{（平板）}
\label{eq:ch10-eta}
\end{equation}

\tocheck{未对照原页面图像：式 (10-115) 首因子在 OCR 中只剩一个 $S$；本节按球对称
Fick 方程的标准解写成 $4\pi R\varphi_SD_ec_S(1/\tanh\varphi_S-1/\varphi_S)$，该式对
$\varphi_S$ 取 $0.1$ 与 $10$ 时分别给出 $\eta\to1$ 与 $\eta\to3/\varphi_S$，
与式~\eqref{eq:ch10-eta} 自洽，但原页面的确切写法未核。}

Thiele 模数的物理意义可由式~\eqref{eq:ch10-thiele} 平方后看出：对球形颗粒
$\varphi_S^2=R^2S_vk_S/D_e=\bigl(RS_vk_Sc^m\bigr)\big/\bigl(D_ec/R\bigr)
=\text{反应速率}/\text{扩散速率}$\src{S2 §10.4.2，书内 p.774 / PDF p.263}。
它反映反应速率与扩散速率的相对大小，也就是内扩散的严重程度。对一级不可逆反应：
颗粒小（$R$ 或 $L$ 小）、活性低（$k_S$ 小）或 $D_e$ 大时 $\varphi$ 小而 $\eta\to1$；
反之 $\varphi$ 越大 $\eta$ 越小。三个区间的操作含义是
\src{S2 §10.4.2，书内 p.774 / PDF p.263}：$\varphi<0.4$ 时 $\eta\approx1$，
\textbf{内扩散可忽略}；$0.4<\varphi<3$ 时 $\eta<1$，内扩散有明显作用；
$\varphi>3$ 时 $\eta\approx1/\varphi$，\textbf{内扩散严重}。因此要消除内扩散
（使 $\eta\approx1$），只有降低 $\varphi$ 一条路：提高 $D_e$（改孔结构）和减小粒度。

\begin{figure}[htbp]
\centering
\begin{tikzpicture}
\begin{axis}[
  width=\onecolwidth,
  height=6.0cm,
  xlabel={内扩散模数 $\varphi$},
  ylabel={有效因子 $\eta$},
  xmode=log, ymode=log,
  xmin=0.1, xmax=100, ymin=0.008, ymax=1.5,
  legend style={at={(0.5,-0.26)}, anchor=north, legend columns=2, draw=none, fill=none},
]
\addplot[mainline, mark=*, mark size=1.2pt] coordinates {
  (0.10,0.9993) (0.15,0.9985) (0.20,0.9973) (0.30,0.9941) (0.40,0.9895)
  (0.60,0.9768) (0.80,0.9598) (1.00,0.9391) (1.50,0.8762) (2.00,0.8060)
  (3.00,0.6716) (4.00,0.5630) (6.00,0.4167) (10.00,0.2700) (20.00,0.1425)
  (50.00,0.0588) (100.00,0.0297)};
\addlegendentry{球形，式~\eqref{eq:ch10-eta}}
\addplot[thinline, dashed, mark=square*, mark size=1.2pt] coordinates {
  (0.10,0.9967) (0.15,0.9926) (0.20,0.9869) (0.30,0.9710) (0.40,0.9499)
  (0.60,0.8951) (0.80,0.8300) (1.00,0.7616) (1.50,0.6034) (2.00,0.4820)
  (3.00,0.3317) (4.00,0.2498) (6.00,0.1667) (10.00,0.1000) (20.00,0.0500)
  (50.00,0.0200) (100.00,0.0100)};
\addlegendentry{平板，式~\eqref{eq:ch10-eta}}
\addplot[accentsoft, dotted, thick, forget plot] coordinates {(2,1.5)(100,0.03)};
\addplot[accentsoft, dotted, thick, forget plot] coordinates {(0.667,1.5)(100,0.01)};
\addplot[accent, densely dotted, forget plot] coordinates {(0.4,0.008)(0.4,1.5)};
\addplot[accent, densely dotted, forget plot] coordinates {(3,0.008)(3,1.5)};
\end{axis}
\end{tikzpicture}
\caption{有效因子随 Thiele 模数的变化（一级反应，双对数坐标）。读图看三处：$\varphi<0.4$
时两曲线都在 $0.95$ 以上（内扩散可忽略）；$\varphi>3$ 后两曲线都变成直线且斜率 $-1$，
即 $\eta\approx1/\varphi$（强内扩散）；两条点线是渐近线 $3/\varphi$（球）与 $1/\varphi$
（板），两条竖虚线为 $\varphi=0.4$ 与 $\varphi=3$。条件与局限：等温、一级不可逆反应、
$D_e$ 为常数、球与平板两种理想几何；$\eta$ 由式~\eqref{eq:ch10-eta} 计算，属模型计算。
\modelmark}
\label{fig:ch10b-etaphi}
\end{figure}

\begin{obs}
读图~\ref{fig:ch10b-etaphi}：$\varphi=0.4$ 处球形 $\eta=0.990$、平板 $\eta=0.950$；
$\varphi=1$ 处两者分别为 $0.939$ 与 $0.762$；$\varphi=10$ 处分别为 $0.270$ 与 $0.100$，
而渐近线 $3/\varphi$ 与 $1/\varphi$ 在同一点给出 $0.300$ 与 $0.100$。$\varphi$ 从 $0.4$
增到 $10$（$25$ 倍），球形 $\eta$ 从 $0.99$ 掉到 $0.27$（$3.7$ 倍）。
\end{obs}

\begin{infer}
（1）$\varphi$ 的 $25$ 倍变化只换来速率 $3.7$ 倍的下降——因为强扩散区里
$\eta\propto1/\varphi$，而 $\varphi\propto R$，所以\textbf{速率大致与粒径成反比}：
粒径减半，总包速率约翻倍，而\textbf{本征速率常数一点没变}。这正是「改变粒径看转化率」
能判别内扩散的原因。（2）平板与球的曲线在两端靠近、在中间分开，说明\textbf{几何形状
只影响系数，不影响标度}；因此工程上常用一个统一长度把不同形状折到同一条曲线上。
（3）$\varphi>3$ 后两条曲线平行（斜率同为 $-1$），故\textbf{强内扩散下表观活化能减半}
——这是内扩散最隐蔽的后果：速率数据看起来仍然「服从 Arrhenius 定律」，只是斜率错了。
\end{infer}

\textbf{表观级数与表观活化能会被改成什么。}对 $m$ 级反应，$r=\eta k_SS_vc_S^m$，
两端取对数后对 $\ln c_S$ 求导得表观级数 $m'=\dd\ln\eta/\dd\ln c_S+m$，即 $m'$ 由
$\eta$ 与 $c_S$ 的关系决定。把 $\varphi$ 写成浓度的函数
$\varphi_S=R\sqrt{S_vk_Sc_S^{m-1}/D_e}$（平板同理）再代入 $\eta\approx1/\varphi$，
可得\src{S2 §10.4.2，书内 p.774—775 / PDF p.263—264}：

\begin{itemize}
\item \textbf{分子扩散}：$D_{ej}$ 与总压成反比（$D_e\propto1/p\propto1/c_S$），
故 $m'=m/2$；
\item \textbf{Knudsen 扩散}：$D_{Ke}$ 与压力无关，故 $m'=(m+1)/2$。
\end{itemize}

\textbf{表观活化能}：定义 $E_a=-R\,\dd(\ln\eta k_S)/\dd(1/T)$，本征值
$E=-R\,\dd(\ln k_S)/\dd(1/T)$，则 $E_a/E$ 是 $\eta$ 的函数；在强内扩散区
（$\varphi>3$、$\eta\approx1/\varphi$）有 $E_a/E\approx1/2$
\src{S2 §10.4.2，书内 p.775 / PDF p.264}。注意这两个结论的\textbf{成立前提}：
$m'$ 的两种结果分别假定了分子扩散与 Knudsen 扩散的\textbf{压力依赖}（并默认理想气体
$c\propto p$，过渡区两者叠加时两个公式都不适用）；$E_a\approx E/2$ 假定了
$D_e$ 与温度无关。\textbf{若把扩散系数的温度依赖算进去，强内扩散的表观活化能应写成
$E_a\approx(E+E_D)/2$}（$E_D$ 为扩散活化能，分子扩散约 $1.5RT$、Knudsen 扩散约
$0.5RT$）——本项为学科标准结果，原书未展开，列此以免把 $E/2$ 当成严格值。

\textbf{不同几何的互检（一个可验算的旁证）。}为使各几何模型可比，引入长度参数
$d=$颗粒体积$/$扩散面积\src{S2 §10.4.3，书内 p.776 / PDF p.265}：平板 $d=L$，
圆球 $d=R/3$（因 $V_p/S_{\text{外}}=R/3$）。于是用统一长度表达的模数为
$\varphi_S/3$——\tocheck{未对照原页面图像：原书此处 OCR 作「以统一长度表达的
$\varphi=\tfrac12\varphi_S$」。按 $\varphi_S=3\varphi_{\text{统一}}$（即 $\varphi_S/3$）
计算：统一模数为 $1$ 时 $\eta_{\text{板}}=0.762$、$\eta_{\text{球}}(\varphi_S=3)=0.672$，
差 $0.090$，恰与原书「对一级反应 $\eta_L$ 最大而 $\eta_S$ 最小，两者相差约 $0.09$」
相符（差值在统一模数 $1.15$ 附近达到极大 $0.092$）；若按 $\tfrac12$ 折算则得 $0.32$，
明显不符。故疑 OCR 的「$1/2$」应为「$1/3$」，引用前请核对页面。}

\subsubsection{复杂情形：双组分、摩尔数变化、双曲线速率方程、可逆反应}

实际体系很少满足上一小节的四条假定，原书逐项讨论了偏离的后果
\src{S2 §10.4.3，书内 p.776—777 / PDF p.265—266}。

\textbf{（1）双组分反应。}反应对两个物种都是一级、其中一个（如 $\mathrm{B_1}$）受内扩散
控制时，以平板模型为例，其内扩散模数为

\begin{equation}
\varphi_{L,1}=L\sqrt{\frac{k_SS_vc_2}{(D_e)_1}}
\label{eq:ch10-bicomp}
\end{equation}

当 $\eta$ 小于 $0.4$ 时它可近似为 $1/\varphi_{L,1}$，此时\textbf{表观级数对 $\mathrm{B_1}$
仍是一级，对 $\mathrm{B_2}$ 却变成 $3/2$ 级}（指 Knudsen 扩散区）。
\tocheck{未对照原页面图像：式 (10-122) 的分子（$c_2$ 是否带幂次、下标 $(D_e)_1$ 的位置）
与「$\eta$ 小于 $0.4$」这一临界表述均按 OCR 转写，数值边界请核页面。}

\textbf{（2）反应伴有摩尔数变化。}在分子扩散或过渡区中，反应总摩尔数增加使反应物更难
扩散进入颗粒，$\eta$ 更低；总摩尔数减少则相反。定量说明需引入摩尔数增加 $\delta$ 与
摩尔分数 $Y_1$ 的乘积作参数：对圆球状催化剂，若 $\varphi_S<3$ 而主要反应物
$\mathrm{B_1}$ 的 $\delta Y_1$ 在 $\pm0.5$ 以内，$\eta$ 的变化不超过 $10\%$。

\textbf{（3）颗粒形状。}除球与板之外还有圆柱、长杆模型。$\eta$—$\varphi$ 曲线对不同
几何并不相同：$\varphi$ 很小或很大时各形状差别相当小，在统一模数 $1$ 附近差别最大
（数值互检见上一小节）；圆柱模型居中。

\textbf{（4）双曲线式速率方程的有效因子。}用 L-H 表达式时 $\eta$ 与 $\varphi$ 的关系
复杂得多，需引入两个组合参数 $K_d$ 与 $v_d$\src{S2 §10.4.3，书内 p.776—777 /
PDF p.265—266}：

\begin{equation}
K_d=\frac{K_1-D_1\sum_i\nu_iK_i/D_i}
{1+\sum_iK_i\bigl(p_i+\nu_ip_1D_1/D_i\bigr)},\qquad
v_d=\frac{D_2p_2}{\nu_2D_1p_1}-1
\label{eq:ch10-kd}
\end{equation}

式中 $K_1$、$D_1$、$p_1$ 是双组分中主要反应物的吸附平衡常数、扩散系数与分压，带下标
$i$ 的是其他组分。选择反应物时要使 $v_d\ge0$。若各 $K_1$ 都很小，$K_d\to0$，反应
趋于一级；对零级反应 $K_dp_1\to\infty$，对高级次反应 $K_dp_1$ 为负；
$K_dp_1$ 最小为 $-1$，对应强产物阻抑。$p_2$ 或 $D_2$ 远大于 $\mathrm{B_1}$ 的相应值时
$v_d$ 较大。原书的 $\eta$—$\varphi_L$ 曲线族（图 10-6）给出三条要点
\src{S2 §10.4.3，书内 p.777 / PDF p.266}：$K_dp_1$ 或 $v_d$ 越大曲线位置越高；
在 $K_dp_1$ 很小时 $v_d$ 的影响很小；\textbf{当 $v_d>0.1$ 且 $K_dp_1\ge10$ 时，
$\eta$ 在某些 $\varphi_L$ 区间内可以比 $1$ 大很多}，而且部分曲线中段的点线连接处
实际是 $\eta$ 多值区间，对应非稳态操作。这两点与\textbf{非恒温}扩散控制的速率行为
相似，读图时不要把它们混为一谈。
\tocheck{未对照原页面图像：式 (10-123) 的完整形式（$K_d$ 分子的符号、$v_d$ 的定义与
$K_dp_1$ 的取值范围）在 OCR 中严重失真，此处按可见结构转写，引用前须核对页面。}

\textbf{（5）可逆反应。}存在显著逆向反应时 $\eta$ 的解析更复杂，一般要在同样的
$\varphi_L$ 下得到\textbf{更低}的 $\eta$。原书给出的一组对比：对一级不可逆反应，
$\varphi_L=0.55$ 及 $3.15$ 时 $\eta$ 约为 $1.0$ 与 $0.30$；若为一级可逆反应、
平衡转化率 $50\%$，同样 $\varphi_L$ 下 $\eta$ 降为 $0.80$ 与 $0.16$。对高级数反应
影响更大。\src{S2 §10.4.3，书内 p.777 / PDF p.266}
\textbf{实用含义}：在接近平衡的条件下用恒温不可逆模型的 $\eta$ 去校正速率，会\textbf{高估}
有效因子，从而低估内扩散的严重程度。

\subsubsection{非恒温粒内传质：非恒温有效因子}

催化剂本身导热率低、扩散系数大、反应热效应大时，颗粒内温度会显著不同于外表面温度。
有效因子仍定义为实测速率与「无粒内传热传质限制」速率的比值，但解热量与物料衡算的
联立方程时须引入两个无因次参数\src{S2 §10.4.4，书内 p.778 / PDF p.267}：

\begin{equation}
\gamma=\frac{E}{RT},\qquad
\beta=\frac{c_s(-\Delta H)D_e}{\lambda_pT_s}
\label{eq:ch10-gammabeta}
\end{equation}

式中 $\Delta H$ 为反应焓变，$\lambda_p$ 为多孔颗粒的热导率。对放热反应 $\beta>0$，其
物理意义是\textbf{稳态下颗粒内可能的最大温差与外表面温度之比}
$(T-T_s)_{\max}/T_s$。原书给出的典型值：$\gamma=20$ 对应 $E\approx100$\,\unit{kJ.mol^{-1}}、
$T\approx600$\,\unit{K}（按 $\gamma=E/RT$ 核算为 $99.8$\,\unit{kJ.mol^{-1}}，自洽），
$\beta$ 在 $\pm0.8$ 之间作出一族 $\eta$—$\varphi_S$ 曲线，$\beta=0$ 即恒温情况。
\src{S2 §10.4.4，书内 p.778 / PDF p.267}

原书对图 10-7 的读法有三条结论，都是恒温模型给不出的
\src{S2 §10.4.4，书内 p.779 / PDF p.268}：（1）$\varphi_S$ 高时各曲线的 $\eta$ 仍与
$\varphi_S$ 成反比，与恒温情况相同；（2）$\beta>0$ 时，\textbf{对一定的 $\varphi_S$，
$\eta$ 可以超过、甚至远超过 $1$}——因为颗粒内温升引起的速率增加远超过浓度降低引起的
速率下降；（3）$\varphi_S$ 较小的区间里 $\eta$ 是\textbf{多值}的，同一个 $\varphi_S$
下可以有三个 $\eta$ 值而放热与散热都能平衡，其中间值对应\textbf{亚稳态}：稍有扰动即
破坏，颗粒温度不是升到与较高 $\eta$ 相应的稳态（控速步骤为气固传质），就是降到与较低
$\eta$ 相应的稳态（表面化学反应控速）。

\begin{pitfall}
\textbf{非恒温体系的两个反直觉之处。}（1）看到 $\eta>1$ 不要先怀疑算错——在放热反应
且粒内温差显著时它是物理上可能的；反过来，\textbf{若用恒温模型解释 $\eta>1$，会把
「本征速率被低温区低估」误读成「模型形式错误」。}（2）$\eta$ 多值意味着\textbf{同一套
操作条件可能出现两个稳定态}，实验上表现为升温和降温测得的速率不重合（滞后）。此时
不能用单值 $\eta$ 曲线做数据校正，必须先判定处在哪个态。
\end{pitfall}

\subsubsection{判据：假设、失效条件与可操作的实验}

判据的作用是把「传递有没有干扰」变成一个\textbf{可算的数}，并用三组实验去检验它。
\textbf{三个判据的前提是同一条}：式~\eqref{eq:ch10-basis} 的速率基准必须统一。

\begin{principle}
\textbf{三个判据（全部写成体积基准；质量基准为等价写法）。}
\textbf{（1）外扩散：Damk\"ohler 数判据}\src{S2 §10.5.4，书内 p.791—792 /
PDF p.280—281}，对 $m$ 级反应
\begin{equation}
Da=\frac{r_vR}{k_mc_B}<\frac{0.15}{m},
\qquad\text{质量基准：}\ Da=\frac{\rho_pr_mR}{k_mc_B}
\label{eq:ch10-da}
\end{equation}
式中 $c_B$ 为组分 $\mathrm{B_1}$ 在流体相的浓度，$R$ 为颗粒半径，$k_m$ 为颗粒外膜的
传质系数。\textbf{量纲}：
$\unit{mol.m^{-3}.s^{-1}}\times\unit{m}\big/\bigl(\unit{m.s^{-1}}\times
\unit{mol.m^{-3}}\bigr)$——无量纲；质量基准式中的
$\unit{kg.m^{-3}}\times\unit{mol.kg^{-1}.s^{-1}}$ 恰好还原为
$\unit{mol.m^{-3}.s^{-1}}$，故两式\textbf{数值恒等}。

\textbf{（2）内扩散：Weisz--Prater 判据}\src{S2 §10.5.4，书内 p.792 / PDF p.281}。
它可以直接由有效因子的定义推出，不必知道本征动力学：把观测速率
$r_v=\eta k_vc_S$ 代入 $\eta\varphi^2$ 得
\begin{equation}
\Phi\equiv\eta\varphi^{2}=\frac{r_vR^{2}}{D_ec_B},
\qquad\text{质量基准：}\ \Phi=\frac{\rho_pr_mR^{2}}{D_ec_B}
\label{eq:ch10-wp}
\end{equation}
式中 $D_e$ 为有效扩散系数。\textbf{量纲}：
$\unit{mol.m^{-3}.s^{-1}}\times\unit{m^2}\big/\bigl(\unit{m^2.s^{-1}}\times
\unit{mol.m^{-3}}\bigr)$——无量纲。由于 $\Phi=\eta\varphi^{2}$ 把两个未知量绑在一起，
判据本身只给 $\Phi$，\textbf{阈值必须按「容许多大的 $\eta$ 下降」来定}。对等温球形
颗粒、一级反应，把式~\eqref{eq:ch10-eta} 的 $\eta(\varphi)$ 逐点代入可算得：
$\Phi=0.3$ 对应 $\varphi\approx0.55$、$\eta\approx0.98$；$\Phi=1$ 对应
$\varphi\approx1.03$、$\eta\approx0.94$。因此原书「一级反应 $\Phi<1$ 可保证有效因子
\textbf{小于} $0.95$」的方向应为\textbf{大于}——$\Phi<1$ 保证的是 $\eta>0.94$；
若要求 $\eta>0.95$ 需 $\Phi<0.74$，要求 $\eta>0.99$（即常用界 $\varphi<0.4$）则需
$\Phi<0.16$。平板几何在同一个 $\Phi$ 下 $\eta$ 更低（$\Phi=1$ 时 $\eta\approx0.70$），
所以\textbf{报阈值必须连几何与级数一起报出}；下文按惯例写「一级 $\Phi<1$、二级
$\Phi<0.3$」，其中 $\Phi<1$ 是\textbf{不出现严重内扩散}的设计上限，不等于
「$\eta$ 仍在 $0.95$ 以上」。
\textbf{注意}：质量基准式里的 $\rho_p$ 就是\textbf{把质量速率换算成体积速率的那一个
因子}，它只出现一次。

\textbf{（3）相间传热：Mears 判据}（$\eta>0.95$ 的条件）
\src{S2 §10.5.4，书内 p.791—792 / PDF p.280—281}：
\begin{equation}
\frac{r_v(-\Delta H_r)R}{h_BT_B}<0.15\,\frac{\mathscr{R}T_B}{E}
\label{eq:ch10-mears}
\end{equation}
式中 $h_B$ 为相间传热系数，$T_B$ 为流体温度，$E$ 为活化能；为免同一字母歧义，
\textbf{右侧的气体常量记作 $\mathscr{R}$}，左侧的 $R$ 仍是颗粒半径。
\textbf{量纲}：左侧
$\unit{mol.m^{-3}.s^{-1}}\times\unit{J.mol^{-1}}\times\unit{m}\big/
\bigl(\unit{W.m^{-2}.K^{-1}}\times\unit{K}\bigr)
=\unit{J.m^{-2}.s^{-1}}\big/\unit{W.m^{-2}}$——无量纲（$\unit{W}=\unit{J.s^{-1}}$）；
右侧 $\unit{J.mol^{-1}.K^{-1}}\times\unit{K}\big/\unit{J.mol^{-1}}$——无量纲。
质量基准式左分子写成 $\rho_pr_m(-\Delta H_r)R$，同样数值恒等。
颗粒半径与几何的关系：由 $S_{\text{外}}/V_p=3/R$ 得 $R=3V_p/S_{\text{外}}$，
非球形颗粒按此折等效半径，不要直接拿筛分粒度。
\end{principle}

\begin{figure}[htbp]
\centering
\begin{tikzpicture}
\begin{axis}[
  width=\onecolwidth,
  height=6.0cm,
  xlabel={颗粒直径 $d_p$ / \unit{mm}},
  ylabel={Weisz--Prater 数 $\Phi$},
  xmode=log, ymode=log,
  xmin=0.05, xmax=10, ymin=1e-7, ymax=1e4,
  legend style={at={(0.5,1.02)}, anchor=south, legend columns=3, draw=none, fill=none},
]
\addplot[mainline, domain=0.05:10, samples=60]{0.075*x^2};
\addlegendentry{正确：$r_v$ 基准}
\addplot[thinline, dashed, domain=0.05:10, samples=60]{6.25e-5*x^2};
\addlegendentry{少乘 $\rho_p$}
\addplot[accentsoft, dashdotted, domain=0.05:10, samples=60]{90*x^2};
\addlegendentry{多乘 $\rho_p$}
\addplot[accent, densely dotted, thick, forget plot] coordinates {(0.05,1)(10,1)};
\addplot[accent, densely dotted, thick, forget plot] coordinates {(0.05,0.3)(10,0.3)};
\end{axis}
\end{tikzpicture}
\caption{同一条 Weisz--Prater 判据、同一个观测速率，三种速率基准下 $\Phi$ 随粒径的变化
（双对数，$\Phi\propto d_p^{2}$）。读图看三线的相对位置：正确基准居中；「少乘 $\rho_p$」
那条线偏低 $1/1200$（判据失效），「多乘 $\rho_p$」那条线偏高 $1199$ 倍（误判），两条
错误线彼此相距 $1.44\times10^{6}$ 倍；两条水平点线为阈值 $\Phi=1$（设计上限）与
$\Phi=0.3$（二级反应用；对一级反应它给出 $\eta>0.95$）。
条件与局限：等温球形颗粒、一级反应，取
$r_m=9$\,\unit{mmol.g^{-1}.h^{-1}}、$\rho_p=1200$\,\unit{kg.m^{-3}}、
$D_e=1.0\times10^{-6}$\,\unit{m^2.s^{-1}}、$c_B=10$\,\unit{mol.m^{-3}}
（教学设定数值），曲线为模型计算。\modelmark}
\label{fig:ch10-wp}
\end{figure}

\begin{obs}
读图~\ref{fig:ch10-wp}：正确基准下 $\Phi=0.075\,d_p^{2}$，与 $\Phi=0.3$ 相交于
$d_p=2.0$\,\unit{mm}，与 $\Phi=1$ 相交于 $d_p=3.65$\,\unit{mm}；「多乘 $\rho_p$」那条线
是 $90\,d_p^{2}$，与 $\Phi=1$ 相交于 $d_p=0.105$\,\unit{mm}；「少乘 $\rho_p$」那条线是
$6.25\times10^{-5}d_p^{2}$，在横轴到 $10$\,\unit{mm} 时 $\Phi$ 仅 $6.3\times10^{-3}$，
始终没有触及任何阈值。三条线在任意横坐标上的比值恒为 $1/1200$（少乘线/正确线）与
$1200$（多乘线/正确线）。三条线的系数可以把量纲检查直接算出来：$r_v=\rho_pr_m
=1200\times2.5\times10^{-3}=3.0$\,\unit{mol.m^{-3}.s^{-1}}，
$\Phi=r_vR^{2}/(D_ec_B)$ 中 $R=d_p/2$，故
$\Phi=r_vd_p^{2}/(4D_ec_B)=3d_p^{2}/(4\times10^{-5})=0.075\,d_p^{2}$
（$d_p$ 以 \unit{mm} 计）；少乘线用 $r_m$ 代替 $r_v$，系数即
$2.5\times10^{-3}/(4\times10^{-5})=6.25\times10^{-5}$；多乘线用 $r_v\rho_p$，
系数即 $0.075\times1200=90$。
\end{obs}

\begin{infer}
（1）三条线来自\textbf{同一组速率数据}，差别只有一个颗粒密度 $\rho_p$；结论却完全相反：
正确基准说「粒径超过约 $2$\,\unit{mm} 就要担心内扩散」，多乘 $\rho_p$ 说「$0.1$\,\unit{mm}
以上都不合格」（会把合格数据误判成受内扩散控制而白白丢弃），少乘 $\rho_p$ 说
「$10$\,\unit{mm} 也通过」（判据失去意义，把扩散控制当本征动力学）。（2）因此
\textbf{判据的第一步不是代入数字，而是写出基准}；判据不自检量纲就等于没做。
（3）判据的数量级特征是 $\Phi\propto d_p^2$：粒径减半，$\Phi$ 降到 $1/4$——这就是
「改变粒径看转化率」这组实验之所以有效的根据。
\end{infer}

\begin{measure}
\textbf{把判据做成实验：三组自变量、三条判读规则。}
\src{S2 §10.5.4，书内 p.791—792 / PDF p.280—281}
\textbf{实验一（外扩散）——同一接触时间下改变线速。}在同一反应器内随质量流速改变而
相应增减催化剂用量，使接触时间 $W/F$ 不变；若质量流速增大而转化率趋于稳定值，
外扩散已消除。也可用截面积相同的两个反应器装不同量催化剂、按比例引入流量：高流速区
两者转化率无差别，低流速区不同。还可在同一反应器中改变流量、在不同温度下求表观速率
常数：不同流量的曲线重合则无外扩散，高温下不重合则外扩散在起作用。\textbf{判读}：
取「流速高到对转化率无明显影响」的区间做动力学实验，再用
式~\eqref{eq:ch10-da} 复核。
\textbf{实验二（内扩散）——同一空速下改变粒径。}从同一批催化剂取不同粒径的几份样品，
测同一空速下的转化率；粒径减小时转化率上升到一极限值，取与该极限值相应的粒径
\textbf{以下}的催化剂作研究，内扩散即可排除。\textbf{判读}：用
式~\eqref{eq:ch10-wp} 复核，且 $\Phi$ 必须随 $d_p^2$ 变化，若实测 $\Phi$ 与粒径无关，
说明 $r_v$、$D_e$ 或 $c_B$ 的取法有问题。
\textbf{实验三（温度）——把三个区画出来。}前两组实验必须在\textbf{整个实验温度区间}
都通过；只在高温区通过、低温区未检验的数据不可外推。\textbf{还要满足理想活塞流的
三个几何条件}：反应管内径至少为颗粒直径的 $10$ 倍以上；气-固反应床层长度至少为粒径的
$50$ 倍以上；液-固反应床层长度至少为粒径的几百倍以上（否则沟流或轴向弥散破坏活塞流
假定）。
\end{measure}

\begin{table}[htbp]
\centering
\caption{三个判据的假设与失效条件。左列三式均用体积基准 $r_v$；换成质量基准时
\textbf{每式恰好乘一个} $\rho_p$——不是每式各自再补一个。
\src{据 S2 §10.5.4，书内 p.790—793 / PDF p.279—282 归纳。}}
\label{tab:ch10-criteria}
\begin{tabularx}{\textwidth}{@{}l X X@{}}
\toprule
判据 & 表达式（体积基准） & 主要假设与失效条件 \\
\midrule
Damk\"ohler 数（外扩散） & $Da=r_vR/(k_mc_B)<0.15/m$ &
假设等温、稳态、外膜传质为串联阻力、对 $\mathrm{B_1}$ 为 $m$ 级。
\textbf{失效}：$m\le0$ 时阈值发散（零级反应不适用）；强放热时必须同时用 Mears 判据；
$k_m$ 须由 $Re$、$Sc$ 在关联式适用范围内取得；非球形颗粒的 $R$ 要按 $3V_p/S_{\text{外}}$
折算 \\
Weisz--Prater 数（内扩散） & $\Phi=r_vR^{2}/(D_ec_B)$，一级 $<1$、二级 $<0.3$ &
假设等温、$D_e$ 为常数、单一反应、几何为球或板、级数已知。
\textbf{失效}：非等温（$\beta$ 不可忽略时 $\eta$ 可大于 $1$，恒温判据给出的「安全」
结论不可靠）；过渡区 $D_e$ 随组成变化；可逆反应与强产物阻抑（同一 $\varphi$ 下 $\eta$
更低甚至多值）；只能判「是否受限制」，不能给出本征参数 \\
Mears 判据（相间传热） & $r_v(-\Delta H_r)R/(h_BT_B)<0.15\mathscr{R}T_B/E$ &
假设稳态、相间膜传热、目标为 $\eta>0.95$、活化能 $E$ 已知、温度用绝对温标。
\textbf{失效}：$E$ 未知时右侧无法计算（用表观 $E$ 会引入偏差）；粒内温差显著时须改用
非恒温有效因子（式~\eqref{eq:ch10-gammabeta} 的 $\beta$）判断，Mears 判据管不到粒内 \\
\bottomrule
\end{tabularx}
\end{table}

\begin{dataflow}
\textbf{一次完整的传递影响判定（判据 + 实验互为回查）。}
（1）写出速率基准：把分析数据换成 $r_m$，再由式~\eqref{eq:ch10-basis} 得 $r_v$；
（2）取几何参数：$R=3V_p/S_{\text{外}}$、$\rho_p$、$c_B$ 与温度；
（3）算 $D_e$：按表~\ref{tab:ch10-diffmech} 判断机制（孔径与平均自由程比较），
分子与 Knudsen 并存时用 $1/D_e=1/D_{Ke}+1/D_{ej}$；
（4）按表~\ref{tab:ch10-criteria} 逐项算 $Da$、$\Phi$，并算热判据；
（5）用实验一、二、三检验判据的结论：判据说「通过」而实验显示转化率随线速或粒径变化，
则优先相信实验——通常是 $D_e$ 取错（无反应下测的值与反应下可差 $4$—$5$ 倍）
或 $k_m$ 的关联式超出适用范围；
（6）判据通过后，才可以在本征动力学区拟合参数，并在论文里同时报出判据的数值与基准。
\end{dataflow}

\begin{pitfall}
\textbf{本节的五个误判，前三个直接毁掉结论。}
（1）\textbf{把扩散控制当本征动力学}：$\varphi>3$ 时速率与粒径成反比、表观活化能减半、
级数被改写（分子扩散 $m/2$、Knudsen $(m+1)/2$），此时得到的「动力学参数」是设备参数
的函数，外推到工业粒度必然出错。识别特征是\textbf{改变线速或粒径时转化率变}。
（2）\textbf{速率基准混用}：把 $\unit{mol.g^{-1}.s^{-1}}$ 的数值直接代入
式~\eqref{eq:ch10-da} 或式~\eqref{eq:ch10-mears}，得到的量单位是
$\unit{m^3.kg^{-1}}$——\textbf{有量纲，无论数值大小都不成判据}；这是「看起来通过判据」
的假阴性来源。\textbf{自检只能靠逐式展开 SI 单位}。
（3）\textbf{重复计入密度}：式~\eqref{eq:ch10-wp} 用 $r_v$ 时不再乘 $\rho_p$，
用 $r_m$ 时乘\textbf{一次} $\rho_p$；写成 $r_v\rho_pR^2/(D_ec_B)$ 等于多算一个
$\rho_p$（图~\ref{fig:ch10-wp} 中就是偏上 $1200$ 倍那条线），而漏乘则偏下 $1200$ 倍。
（4）\textbf{忽略热点与温度梯度}：$\eta>1$、同一操作点两个稳定态、升降温数据不重合，
都可能是非恒温的表现，而不是「模型错了」；强放热反应里相间温差可远大于粒内温差。
（5）\textbf{用 $R^2$ 判定模型优劣}：判据与模型检验的判据是残差结构与参数稳定性，
$R^2$ 高只说明曲线穿过了数据点。
\end{pitfall}

\begin{crosscheck}
\textbf{判据之间、判据与实验之间的回查关系。}
（1）\textbf{顺序}：外扩散与内扩散是串连的（外扩散在前），必须先排外扩散再排内扩散；
若 $Da$ 不通过，$\Phi$ 里用的 $r_v$ 本身就含外扩散的影响，判据失去意义。
（2）\textbf{两类判据互补}：$Da$ 与 $\Phi$ 管等温情形，Mears 判据管相间传热，粒内温差
由 Prater 数 $\beta$（式~\eqref{eq:ch10-gammabeta}）管——三者的检查对象不同，
不能用一个代替另一个。
（3）\textbf{与 $D_e$ 的测量方式回查}：无反应下估算或测定的扩散系数与反应条件下的
实际有效扩散系数\textbf{可差 $4$—$5$ 倍}（反应与非反应时扩散通量比相差甚大，且与表面
扩散、孔径分布不均有关），因此求内扩散参数时优先用\textbf{反应条件下}的数据
（改变粒径求有效因子、单锭扩散反应器等），具体方法见本章后半部分。
（4）\textbf{与独立实验回查}：判据说「本征区」而瞬态或同位素实验给出的吸附量与由稳态
数据反推的 $K$ 矛盾时，问题通常不在判据，而在数据窗口没有覆盖机理发生变化的区间。
\end{crosscheck}
